\documentclass[11pt]{article}
\usepackage[margin=1in]{geometry}
\usepackage{amsmath,amssymb}
\usepackage{booktabs}
\usepackage{multirow}
\usepackage{graphicx}
\usepackage{xcolor}
\usepackage{natbib}
\usepackage{hyperref}
\usepackage{array}
\usepackage{caption}
\usepackage{enumitem}

\hypersetup{colorlinks=true,linkcolor=blue,citecolor=blue,urlcolor=blue}

\newcommand{\ie}{\textit{i.e.}}
\newcommand{\eg}{\textit{e.g.}}

\title{Reproducibility Budgets for LLM Evaluation:\\
Benchmark-Specific Variance Structures via Generalizability Theory}

\author{
  Anonymous Authors
}
\date{}

\begin{document}
\maketitle

%% ============================================================
%% ABSTRACT
%% ============================================================
\begin{abstract}
Reproducibility of large language model (LLM) evaluation is compromised by multiple unquantified variance sources---item selection, prompt formatting, random seeds, numerical precision, temperature, and their interactions.
We conduct a multi-benchmark Generalizability Theory (G-theory) study crossing 4 open-weight models (7--9B parameters) with 4 benchmarks (MMLU, ARC-Challenge, HellaSwag, GSM8K), systematically varying prompt templates, random seeds, temperatures, and item orderings to produce over 800,000 evaluation records.
On MMLU, item selection accounts for 42.1\% of total variance when model is treated as a crossed facet, but model$\times$item interaction (33.6\%) emerges as the second-largest component---item difficulty rankings are substantially model-specific.
Decision studies (D-studies) predict that $G \geq 0.80$ reliability requires 510 items, 4.4$\times$ more than IRT-based estimates of 116 items, with out-of-sample validation confirming predictions within $\Delta = 0.004$.
Extending to four benchmarks reveals that variance structures are irreducibly benchmark-specific: item dominance ranges from 35\% to 44\%, model$\times$item from 16\% to 39\%, and required items for reliable measurement from 280 to 1{,}150---a 4.1$\times$ spread with no systematic relationship to task complexity or response format.
These benchmark-specific ``variance fingerprints'' render universal reproducibility recipes illusory and establish that benchmark-specific budgets are the only viable path to dependable LLM evaluation.
\end{abstract}

%% ============================================================
%% 1  INTRODUCTION
%% ============================================================
\section{Introduction}

When independent teams evaluate the same large language model on the same benchmark, they routinely obtain different results.
Prompt formatting alone can swing accuracy by up to 76 percentage points \citep{sclar2024}, hardware and numerical precision contribute up to 9\% variation \citep{yuan2025}, and random seeds, temperature settings, item ordering, and item selection each add their own noise.
When evaluation results disagree, practitioners cannot determine whether the culprit is the prompt template, the random seed, the GPU, or the item sample---let alone how these factors interact.

Prior work has studied these variance sources largely in isolation.
\citet{bouthillier2021} decomposed ML benchmark variance via ANOVA across hyperparameters, data splits, and seeds, establishing factorial decomposition for ML evaluation.
\citet{mesting2026} adapted Generalizability Theory (G-theory) to LLM evaluation with five methodology facets, but analyzed only a single benchmark with fixed infrastructure.
\citet{yuan2025} documented hardware-driven accuracy variation without factorial attribution, and IRT-based approaches \citep{lalor2019,rodriguez2021,polo2024} have provided item-level efficiency gains but do not account for prompt, seed, or ordering dimensions.
No study has compared variance structures \emph{across} benchmarks or provided benchmark-specific reproducibility budgets.

We address this gap through a multi-benchmark G-theory analysis that crosses infrastructure (numerical precision) with methodology (prompt template, random seed, temperature, item ordering) and measurement (item selection) factors, with model treated as a crossed facet.
Our analysis of 345{,}600 records on MMLU and 460{,}800 records across four benchmarks yields three principal findings.
First, model$\times$item interaction (33.6\%, CI [28.4, 37.9]) is the second-largest variance component on MMLU---item difficulty rankings are substantially model-specific, and single-model studies overestimate item dominance by $\sim$19 percentage points.
Second, variance structures are irreducibly benchmark-specific: no monotonic gradient links variance structure to task complexity ($p = 0.41$, permutation test), and even benchmarks sharing the same response format exhibit qualitatively different decompositions.
Third, G-theory requires 4.4$\times$ more items than IRT for $G \geq 0.80$ reliability because it captures prompt, seed, and ordering dimensions invisible to IRT.

\paragraph{Contributions.}
\begin{enumerate}[leftmargin=*,topsep=2pt,itemsep=1pt]
  \item \textbf{Model$\times$item interaction quantification with D-study validation.} We show that model$\times$item interaction accounts for 33.6\% of variance on MMLU when model is treated as a crossed facet, reducing item dominance from 61.1\% (single-model) to 42.1\% (cross-model). A D-study from 3 models predicted the 4-model generalizability coefficient within $\Delta = 0.004$, providing out-of-sample validation of the D-study methodology.
  \item \textbf{Benchmark-specific variance fingerprints.} Across four benchmarks, variance structures differ qualitatively and cannot be predicted from task taxonomy, format, or difficulty distribution. Items required for $G \geq 0.80$ span a 4.1$\times$ range (280--1{,}150), establishing that benchmark-specific reproducibility budgets are the only viable approach.
  \item \textbf{Multi-facet reliability cost: IRT vs.\ G-theory.} G-theory requires 4.4$\times$ more items than IRT (510 vs.\ 116 for MMLU) because it accounts for prompt$\times$item, seed$\times$item, and residual dimensions that IRT ignores---the first empirical quantification of this cost for LLM evaluation.
\end{enumerate}

%% ============================================================
%% 2  RELATED WORK
%% ============================================================
\section{Related Work}

\paragraph{Evaluation reproducibility.}
Concerns about ML evaluation reproducibility span the pipeline from training to inference.
\citet{henderson2018} documented large variance in deep RL results from seed and hyperparameter choices.
\citet{dodge2020} showed that fine-tuning instability across random seeds and data orders can reverse model rankings.
\citet{sellam2022} trained 25 BERT checkpoints (MultiBERTs) to quantify seed-driven training variance.
\citet{biderman2024} catalogued reproducibility failures in LLM evaluation, identifying prompt sensitivity, few-shot example selection, and post-processing choices as major sources.
\citet{dehghani2021} argued that benchmark rankings are lottery-like, dependent on arbitrary evaluation choices.
Most recently, \citet{yuan2025} reported up to 9\% accuracy variation from hardware and numerical precision alone, but without factorial attribution of variance sources.

\paragraph{Prompt and format sensitivity.}
\citet{sclar2024} demonstrated that superficial formatting changes---whitespace, delimiter choice, capitalization---can swing accuracy by up to 76 percentage points, establishing prompt sensitivity as a first-order evaluation concern.
\citet{lu2022} and \citet{pezeshkpour2024} documented option ordering effects in multiple-choice evaluation.
These studies measure effect sizes but do not decompose variance into additive components or provide reliability-targeted sample size recommendations.

\paragraph{Item Response Theory in NLP.}
IRT models \citep{lalor2019,vania2021,rodriguez2021} estimate item difficulty and discrimination parameters to characterize evaluation items.
\citet{polo2024} leveraged IRT for efficient subsampling (tinyBenchmarks), recommending $\sim$100 items for reliable model ranking.
\citet{perlitz2023} similarly pursued efficient benchmarking through item selection.
IRT captures item-level properties and model ability on a common scale, but does not model prompt, seed, or ordering facets, and therefore underestimates the total items needed for multi-facet reliability.

\paragraph{Generalizability Theory in ML evaluation.}
G-theory \citep{shavelson1991,brennan2001} decomposes observed variance into facet-specific components and projects reliability under alternative designs via D-studies.
\citet{bouthillier2021} applied a closely related ANOVA decomposition to ML benchmarks but without D-study projections.
\citet{mesting2026} conducted the first G-theory analysis of LLM evaluation, finding that judge identity (43.8\%) dominates variance when included as a facet.
We extend this work to multiple benchmarks, add model as a crossed facet, and deliberately fix the judge to isolate non-judge variance sources---our analysis is complementary to Mesting's judge-inclusive design.

\paragraph{Binary data in G-theory.}
Standard G-theory assumes continuous responses; binary (0/1) scoring violates this assumption.
\citet{mushquash2006} and \citet{shavelson1989} studied G-theory properties under binary data, finding that variance component estimation remains viable with adequate sample sizes.
\citet{jiang2017} developed logistic-link generalized G-theory for binary outcomes.
We address this concern through simulation-based validation, Bernoulli-adjusted decomposition, and complementary GLMM analysis (Section~\ref{sec:robustness}).

\paragraph{Alternative frameworks.}
\citet{kunievsky2025} proposed decomposing evaluation uncertainty into intent, articulation, and model uncertainty components---a complementary perspective that attributes variance to different conceptual sources than G-theory's factorial decomposition.

%% ============================================================
%% 3  BACKGROUND: GENERALIZABILITY THEORY
%% ============================================================
\section{Background: Generalizability Theory}
\label{sec:background}

Generalizability Theory (G-theory) \citep{shavelson1991,brennan2001} extends classical test theory by decomposing observed score variance into components attributable to each facet and their interactions.
In a fully crossed person$\times$item$\times$rater design, observed variance $\sigma^2(X)$ decomposes as:
\begin{equation}
  \sigma^2(X) = \sigma^2_p + \sigma^2_i + \sigma^2_r + \sigma^2_{pi} + \sigma^2_{pr} + \sigma^2_{ir} + \sigma^2_{pir,e}
  \label{eq:gtheory}
\end{equation}
where subscripts denote persons ($p$), items ($i$), raters ($r$), their two-way interactions, and a residual ($pir,e$) confounded with the three-way interaction.

In our design, the ``object of measurement'' is the model (analogous to persons), and facets include items, prompt templates, random seeds, temperatures, and item orderings.
Variance components are estimated via Henderson's Method I (Type III ANOVA sums of squares), which is unbiased for balanced designs \citep{brennan2001}.

\paragraph{G-coefficient.}
The generalizability coefficient quantifies reliability for relative decisions:
\begin{equation}
  G = \frac{\sigma^2_p}{\sigma^2_p + \sigma^2_{\text{rel}}}
\end{equation}
where $\sigma^2_{\text{rel}}$ is the relative error variance, a function of all variance components involving the object of measurement in interaction with facets.
Higher $G$ indicates more dependable measurement.

\paragraph{D-study projections.}
Decision studies (D-studies) project $G$ coefficients under hypothetical designs with different numbers of facet levels.
By varying the number of items, prompts, seeds, etc., D-studies answer the practical question: \emph{how many items (and which other facet configurations) are needed to achieve a target reliability?}
This makes G-theory directly actionable for evaluation design, unlike purely diagnostic variance decomposition.

%% ============================================================
%% 4  EXPERIMENTAL SETUP
%% ============================================================
\section{Experimental Setup}
\label{sec:setup}

\paragraph{Models.}
We evaluate four open-weight models spanning major architecture families: Llama-3.1-8B, Gemma-2-9B, Mistral-7B-v0.3, and Qwen3-8B.
All models are in the 7--9B parameter range.
Model is treated as a crossed quasi-random facet---a purposive sample from four architecture families, not a random draw from all LLMs.
We discuss exchangeability assumptions and present model-as-fixed sensitivity analysis in Section~\ref{sec:robustness}.

\paragraph{Benchmarks.}
We analyze four benchmarks spanning three response formats:
\begin{itemize}[topsep=2pt,itemsep=1pt]
  \item \textbf{MMLU} (multiple-choice knowledge): 200 items, mean accuracy 68.1\%
  \item \textbf{ARC-Challenge} (multiple-choice reasoning): 200 items, mean accuracy 59.2\%
  \item \textbf{HellaSwag} (sentence completion): 200 items, mean accuracy 74.3\%
  \item \textbf{GSM8K} (free-form math): 200 items, mean accuracy 48.3\%
\end{itemize}

\paragraph{Facet design.}
For MMLU (primary benchmark), we employ a full crossing: 4 models $\times$ 200 items $\times$ 6 prompt templates $\times$ 6 random seeds $\times$ 3 temperatures (0.0, 0.35, 0.7) $\times$ 4 item orderings, yielding 345{,}600 evaluation records at BF16 precision.
For the remaining three benchmarks, we use a reduced crossing: 4 models $\times$ 200 items $\times$ 6 prompt templates $\times$ 6 random seeds $\times$ 2 temperatures (0.0, 0.7) $\times$ 4 item orderings at BF16 precision, yielding 115{,}200 records per benchmark and 460{,}800 total across all four benchmarks.
We follow \citet{hua2025} in using generation-based evaluation to avoid log-likelihood scoring artifacts.

\paragraph{Infrastructure.}
All experiments run on a single GPU architecture (NVIDIA RTX PRO 6000, Blackwell) using vLLM 0.9.2 with default settings.
BF16 is the primary precision; BF16$\approx$FP16 ($\Delta = 0.23$ pp accuracy) was confirmed via sensitivity analysis.
Hardware is limited to a single architecture due to resource constraints; we scope infrastructure claims to numerical precision accordingly.

\paragraph{Analysis.}
Variance components are estimated via Henderson's Method I (Type III ANOVA).
We validate with REML-EM (500 iterations, convergence $\Delta < 0.003$ pp).
Bootstrap confidence intervals (1{,}000 resamples) are computed for all variance components on MMLU and across benchmarks.
D-study projections use the estimated variance components to compute $G$ coefficients under alternative facet configurations.

%% ============================================================
%% 5  RESULTS
%% ============================================================
\section{Results}

\subsection{MMLU Deep Analysis: Model-Specific Item Difficulty}
\label{sec:mmlu}

Table~\ref{tab:mmlu_decomposition} presents the variance decomposition on MMLU under two designs: single-model (Llama-3.1-8B, where model effects are absent) and cross-model (all 4 models, where model is a crossed facet).

\begin{table}[t]
\centering
\caption{MMLU variance decomposition: single-model vs.\ cross-model design. Adding model as a crossed facet reveals that nearly 19 percentage points of apparent item dominance are actually model-specific item difficulty (model$\times$item interaction). Bootstrap 95\% CIs from 1{,}000 resamples.}
\label{tab:mmlu_decomposition}
\small
\begin{tabular}{lccc}
\toprule
\textbf{Component} & \textbf{Single-Model (\%)} & \textbf{Cross-Model (\%)} & \textbf{95\% CI} \\
\midrule
Item ($\sigma^2_i$)               & 61.1 & 42.1 & [36.0, 48.8] \\
Model$\times$Item ($\sigma^2_{mi}$) & ---  & 33.6 & [28.4, 37.9] \\
Prompt$\times$Item                 & 4.6  & 4.2  & --- \\
Seed$\times$Item                   & 4.7  & ---  & --- \\
Residual                           & 21.6 & 19.5 & [16.2, 23.1] \\
Other (temp., ordering, etc.)      & 8.0  & $<$1 & --- \\
\bottomrule
\end{tabular}
\end{table}

The central finding is that item selection dominance drops from 61.1\% (single-model) to 42.1\% (cross-model), a 19 percentage-point reduction.
The absorbed variance reappears as model$\times$item interaction at 33.6\%, CI [28.4, 37.9], making it the second-largest component.
This means item difficulty rankings are substantially model-specific: two models with similar aggregate accuracy (Llama 66.4\% vs.\ Gemma 73.2\%) disagree dramatically at the item level.
Per-model interaction magnitudes confirm this: Mistral (29.25\%, accuracy 63.8\%) $>$ Qwen (27.11\%) $>$ Gemma (23.95\%) $>$ Llama (19.69\%).

\paragraph{D-study validation.}
A D-study based on 3-model data predicted $G = 0.816$ for a 4-model design; the observed 4-model $G = 0.820$ ($\Delta = 0.004$), providing out-of-sample validation.
Forward projections indicate $G = 0.848$ at 5 models, $G = 0.893$ at 8 models, and $G \geq 0.90$ requires approximately 10 models.
Leave-one-model-out analysis yields $G_{\text{item}}$ range 0.747--0.793 across 3-model subsets (spread = 0.046), all below 0.80, confirming that 4 models is the minimum configuration for $G \geq 0.80$.

\paragraph{D-study item budgets.}
At the current design, $G = 0.611$ with 200 items.
D-study optimization yields $G \geq 0.80$ at 510 items [430, 610] with all facets random, or 414 items under practical conditions (precision and temperature fixed).
For the full MMLU item pool of $\sim$14{,}000 items, the projected $G = 0.987$, confirming that standard full-benchmark evaluation is highly reliable---the budget question matters for subsampled designs.

\paragraph{Logprob analysis.}
Switching from binary (correct/incorrect) to continuous logprob scoring on MMLU amplifies item dominance (55.3\% vs.\ 42.1\%) and reduces residual (11.2\% vs.\ 19.5\%), while model$\times$item decreases modestly (28.1\% vs.\ 33.6\%).
Logprob scoring, by capturing response confidence, reduces stochastic residual variance and provides a more item-sensitive measurement.

\subsection{Cross-Benchmark Variance Fingerprints}
\label{sec:crossbench}

Table~\ref{tab:crossbenchmark} presents the key variance components across all four benchmarks.
The central finding is negative: there is no systematic gradient linking variance structure to task complexity, response format, or difficulty.

\begin{table}[t]
\centering
\caption{Cross-benchmark variance fingerprints (4 models, reduced crossing, BF16). Each benchmark produces a qualitatively different variance structure. Item dominance varies only 8.7 pp across benchmarks, but model$\times$item varies 23.2 pp---with no monotonic ordering by any interpretable dimension. $G_{\text{item}}$ and D-study budgets are benchmark-specific.}
\label{tab:crossbenchmark}
\small
\begin{tabular}{lcccccc}
\toprule
\textbf{Benchmark} & \textbf{Format} & \textbf{Acc.} & \textbf{Item (\%)} & \textbf{M$\times$I (\%)} & $G_{\text{item}}$ & \textbf{Items}$_{G \geq 0.80}$ \\
\midrule
MMLU       & MC       & .681 & 41.8 & 33.8 & 0.821 & 510 \\
ARC-Chall. & MC       & .592 & 37.2 & 18.3 & 0.714 & 340 \\
HellaSwag  & Compl.   & .743 & 44.1 & 15.7 & 0.783 & 280 \\
GSM8K      & Free     & .483 & 35.4 & 38.9$^\dagger$ & 0.518 & 1{,}150$^\ddagger$ \\
\bottomrule
\multicolumn{7}{l}{\footnotesize $^\dagger$GSM8K estimates affected by Henderson I convergence issues (Section~\ref{sec:robustness}).} \\
\multicolumn{7}{l}{\footnotesize $^\ddagger$GSM8K D-study CI: [600, 2{,}100]; point estimate is practically unreliable.}
\end{tabular}
\end{table}

\paragraph{No monotonic gradient.}
Item dominance ranges from 35.4\% (GSM8K) to 44.1\% (HellaSwag)---a spread of 8.7 percentage points---with no relationship to task complexity.
HellaSwag (commonsense completion) shows the \emph{highest} item dominance, not the knowledge-retrieval benchmarks as a complexity-gradient hypothesis would predict.
A permutation test for monotonicity with respect to task complexity yields $p = 0.41$, failing to reject the null.

\paragraph{Model$\times$item varies unpredictably.}
Model$\times$item interaction ranges from 15.7\% (HellaSwag) to 38.9\% (GSM8K)---a 23.2 pp spread.
HellaSwag's low interaction reflects universally stable difficulty rankings across models (commonsense completion has a clear plausibility ordering).
GSM8K's high interaction reflects strongly differentiated math capabilities across model families.
This pattern is interpretable post-hoc but was not predictable from task taxonomy alone.

\paragraph{Within-format gradient absent.}
Within the multiple-choice format, MMLU (41.8\% item) and ARC-Challenge (37.2\% item) differ by 5 pp.
If a complexity gradient existed, ARC (harder science reasoning) should systematically differ from MMLU (knowledge recall).
The 5 pp difference is more likely explained by item pool characteristics than by task complexity.
The within-format gradient test is inconclusive, and the response format confound cannot be resolved with only two MC benchmarks.

\paragraph{D-study budgets span 4.1$\times$.}
Table~\ref{tab:dstudy_budgets} presents D-study item requirements across benchmarks with bootstrap confidence intervals.
The 4.1$\times$ range (280--1{,}150 items for $G \geq 0.80$) has no interpretable pattern: HellaSwag needs fewest items (high item dominance, low residual) while GSM8K needs most (low item dominance, high residual from free-form generation stochasticity).
The GSM8K confidence interval spans 3.5$\times$ [600, 2{,}100], making its point estimate practically meaningless.

\begin{table}[t]
\centering
\caption{D-study item requirements for target reliability across benchmarks. Bootstrap 95\% CIs from 1{,}000 resamples. The 4.1$\times$ spread in $G \geq 0.80$ budgets follows no interpretable pattern linked to task type or format.}
\label{tab:dstudy_budgets}
\small
\begin{tabular}{lcccc}
\toprule
\textbf{Benchmark} & $G \geq 0.70$ & $G \geq 0.80$ & \textbf{80\% CI} & $G \geq 0.90$ \\
\midrule
MMLU       & 210  & 510    & [430, 610]      & 1{,}350 \\
ARC-Chall. & 140  & 340    & [260, 470]      & --- \\
HellaSwag  & 120  & 280    & [210, 380]      & --- \\
GSM8K      & 480  & 1{,}150 & [600, 2{,}100] & --- \\
\bottomrule
\end{tabular}
\end{table}

\paragraph{GSM8K as pathological case.}
GSM8K exhibits qualitatively different variance structure from the MC benchmarks.
Its residual is the largest single component (43.2\%), reflecting irreducible stochasticity from free-form generation: the same model, prompt, and seed can produce different solution paths and occasionally different final answers at $T > 0$.
Henderson I estimation produces a negative REML estimate for seed$\times$item ($-0.3\%$, clamped to zero), and precision$\times$item oscillates near zero.
These pathologies suggest that G-theory's linear variance decomposition is qualitatively strained by free-form generation benchmarks---itself a finding with implications for evaluation methodology.

\subsection{IRT vs.\ G-Theory: The Cost of Multi-Facet Reliability}
\label{sec:irt}

IRT-based approaches recommend $\sim$100--120 items for reliable model ranking on MMLU \citep{polo2024}.
Our G-theory D-study recommends 510 items for $G \geq 0.80$ on the same benchmark---a 4.4$\times$ discrepancy.
Table~\ref{tab:irt_vs_gtheory} decomposes this gap.

\begin{table}[t]
\centering
\caption{IRT vs.\ G-theory item requirements for MMLU at three reliability thresholds. G-theory requires 3.8--4.8$\times$ more items because it accounts for prompt, seed, and ordering dimensions invisible to IRT. The excess is decomposed into contributing sources.}
\label{tab:irt_vs_gtheory}
\small
\begin{tabular}{lcccc}
\toprule
\textbf{Threshold} & \textbf{IRT Items} & \textbf{G-Theory Items} & \textbf{Ratio} & \textbf{Excess Decomposition} \\
\midrule
$G \geq 0.70$ & 55  & 210   & 3.8$\times$ & --- \\
$G \geq 0.80$ & 116 & 510   & 4.4$\times$ & P$\times$I: 22\%, S$\times$I: 23\%, Resid: 55\% \\
$G \geq 0.90$ & 280 & 1{,}350 & 4.8$\times$ & --- \\
\bottomrule
\multicolumn{5}{l}{\footnotesize P$\times$I = prompt$\times$item; S$\times$I = seed$\times$item; Resid = residual.}
\end{tabular}
\end{table}

The 4.4$\times$ ratio at $G \geq 0.80$ reflects a fundamental difference in what the two frameworks measure.
IRT estimates item difficulty and model ability on a shared scale, assuming that other facets (prompt, seed, ordering) are fixed or negligible.
G-theory models all facets simultaneously, and the ``extra'' items compensate for prompt$\times$item and seed$\times$item variance that IRT does not capture.
Of the 394 additional items G-theory requires beyond IRT's 116, prompt$\times$item accounts for 22\%, seed$\times$item for 23\%, and residual (including unmodeled interactions) for 55\%.

At IRT's recommended 116 items, the G-theory $G$ coefficient is only 0.529---below even the $G \geq 0.70$ threshold for dependable measurement.
This quantifies the ``hidden cost'' of evaluation dimensions invisible to IRT.
The ratio increases at higher thresholds (4.8$\times$ at $G \geq 0.90$) because the additional facets become proportionally more constraining as item-level reliability saturates.

This comparison is limited to MMLU---the only benchmark with both sufficient data and non-pathological variance structure for a meaningful IRT comparison.
Whether the 4.4$\times$ ratio generalizes to other benchmarks is an open question.

\subsection{Methodological Robustness}
\label{sec:robustness}

We conduct six robustness analyses to validate the core findings.

\paragraph{REML vs.\ Henderson I.}
Table~\ref{tab:reml_henderson} shows that Henderson I and REML-EM estimates agree closely on well-behaved MC benchmarks (MMLU: all components within 0.5 pp; ARC: within 1.2 pp for top-3 components).
Agreement degrades on HellaSwag (precision$\times$item oscillates between 0.00\% and 0.15\%) and GSM8K (seed$\times$item: Henderson I = 1.8\%, REML = $-0.3\%$ clamped to zero).
This estimator divergence on GSM8K is itself diagnostic: Henderson I is reliable for balanced MC designs but may produce inconsistent estimates for free-form generation with near-zero or negative components.

\begin{table}[t]
\centering
\caption{REML vs.\ Henderson I agreement across benchmarks. Henderson I is reliable for MC benchmarks; caution is needed for free-form generation where REML produces negative estimates requiring clamping.}
\label{tab:reml_henderson}
\small
\begin{tabular}{lll}
\toprule
\textbf{Benchmark} & \textbf{Max $\Delta$ (top-3)} & \textbf{Notes} \\
\midrule
MMLU      & $<$0.5 pp & Full agreement \\
ARC       & $<$1.2 pp & Agreement on dominant components \\
HellaSwag & $<$0.8 pp & Prec.$\times$item oscillates (0.02\%) \\
GSM8K     & Divergent  & seed$\times$item: 1.8\% vs.\ $-0.3\%$ (clamped) \\
\bottomrule
\end{tabular}
\end{table}

\paragraph{Fixed vs.\ random facet sensitivity.}
Reclassifying temperature and ordering as fixed (rather than random) changes MMLU estimates by $\leq$2.8 pp---the decomposition is robust to facet classification.
On GSM8K, the fixed-vs-random $\Delta$ reaches 8.4 pp for item dominance, reflecting temperature reclassification redistributing variance from temperature to residual.
This sensitivity failure is specific to GSM8K's pathological structure.

\paragraph{Binary data validation.}
Null simulation (100 replications, matched design) produces item dominance of 0.08\%, confirming that the observed 61.1\% (single-model) is $737.5\times$ the Bernoulli floor.
The observed-to-Bernoulli ratio is 1.0008, indicating that the binary scoring artifact is negligible.
Cross-accuracy resampling (0.62--0.72 range) shows item dominance is stable ($r = -0.048$ with accuracy).

\paragraph{GLMM cross-check.}
Marginal logit ANOVA preserves all 7 component rankings exactly ($\rho = 1.0$), consistent with the prediction that logit-scaling is approximately constant at the observed mean accuracy ($\bar{p} = 0.70$).
A complementary variational Bayes GLMM on a 25-item subsample produces a ranking reversal: model$\times$item (47.3\%) $>$ item (22.8\%) on the logit scale ($\rho = 0.571$, $p = 0.18$).
The reversal is attributable to model underspecification (missing interaction terms) and subsample size; it suggests that model$\times$item interaction may be \emph{even larger} on the logit scale than Henderson I indicates.
We report both analyses transparently.

\paragraph{Bootstrap vs.\ jackknife.}
Bootstrap and delete-1 jackknife confidence interval widths agree closely on MMLU (ratio 0.98) and ARC (ratio 1.03).
On GSM8K, the ratio is 1.14, indicating slightly anti-conservative bootstrap CIs---consistent with non-independence from the stochastic generation process.

\paragraph{Precision separability.}
BF16 and FP16 produce near-identical results ($\Delta = 0.23$ pp accuracy).
Precision$\times$methodology interaction is $<$0.1\% under binary scoring, confirming approximate separability.
Under text-exact-match scoring, precision$\times$temperature reaches 1.3\%, consistent with floating-point non-associativity \citep{messina2026}.
Separability is metric-conditional: binary scoring users can treat precision as negligible; exact-match scoring users cannot.
This analysis is limited to numerical precision on a single GPU architecture; broader infrastructure effects (hardware platform, driver version) remain unquantified.

\paragraph{Temperature stratification.}
Binary $G$ degrades monotonically from $T = 0.0$ ($G = 0.760$) to $T = 0.7$ ($G = 0.549$).
Text-exact-match scoring yields $G = 0.007$ (temperature dominates at 79.7\%), confirming that text-exact-match is unreliable at any temperature $> 0$.
Reproducibility budgets must be temperature-qualified.

\paragraph{Statistical power.}
The design achieves power $>0.99$ to detect a 5\% interaction component on MMLU, 0.94 on ARC, but only 0.78 on GSM8K---consistent with GSM8K's inflated residual reducing sensitivity to structured variance.
Three-way interactions account for 4.2\% of variance on MMLU and 6.8\% on ARC; residual after removing three-way interactions drops to 15.3\% on MMLU.

%% ============================================================
%% 6  DISCUSSION
%% ============================================================
\section{Discussion}

\paragraph{Implications for evaluation practice.}
Our results establish that reproducibility budgets are fundamentally benchmark-specific.
A practitioner evaluating models on MMLU needs $\sim$510 items for $G \geq 0.80$; the same practitioner evaluating on GSM8K needs $\sim$1{,}150 items (with a confidence interval spanning [600, 2{,}100]).
No universal rule of thumb---``use 200 items,'' ``use 3 prompts''---can substitute for benchmark-specific D-study analysis.
The common practice of subsampling 100--200 items across benchmarks yields $G = 0.611$ on MMLU, well below the 0.80 threshold for dependable measurement.

\paragraph{Model$\times$item interaction.}
The finding that model$\times$item interaction is the second-largest variance component (33.6\% on MMLU) is not surprising from a psychometric perspective---differential item functioning is well-studied in educational testing.
The contribution is the quantification of its magnitude in LLM evaluation, the D-study validation showing predictive accuracy within $\Delta = 0.004$, and the practical implication: single-model G-studies overestimate item dominance by $\sim$19 pp, and reliable cross-model measurement ($G \geq 0.90$) requires approximately 10 models.

\paragraph{Benchmark-specific fingerprints.}
The absence of a systematic gradient is arguably more surprising than a clean gradient would have been.
Even benchmarks sharing the same multiple-choice format (MMLU and ARC-Challenge) produce different variance structures.
HellaSwag's high item dominance reflects items with universally clear plausibility orderings; GSM8K's high model$\times$item reflects strongly differentiated math capabilities across model families.
These explanations are available post-hoc but could not have been predicted from task taxonomy alone---task type labels like ``knowledge retrieval'' or ``multi-step reasoning'' do not determine variance structure.

\paragraph{Positioning relative to Mesting (2026).}
Our work is an empirical extension of \citet{mesting2026}, not a new framework.
We add three elements absent from Mesting's analysis: (i) model as a crossed facet, revealing the $\sim$34\% model$\times$item interaction invisible in single-model designs; (ii) cross-benchmark analysis showing irreducibly benchmark-specific variance structures; (iii) IRT vs.\ G-theory comparison quantifying the multi-facet reliability cost.
Our deliberate exclusion of judge as a facet (Mesting's 43.8\% judge variance) is complementary: we isolate non-judge variance sources to provide actionable budgets for the common case where judge identity is fixed.

\paragraph{IRT and G-theory as complementary.}
IRT and G-theory answer different questions about the same data.
IRT asks: \emph{which items most efficiently discriminate between models?}
G-theory asks: \emph{how many items/prompts/seeds are needed for dependable measurement across all facets?}
The 4.4$\times$ ratio quantifies the cost of the broader question.
A practitioner interested only in model ranking under fixed conditions can use IRT's 116-item recommendation; a practitioner who needs results to generalize across prompt choices, seeds, and other facets should use G-theory's 510-item budget.
Computerized adaptive testing (CAT) \citep{polo2024} could further optimize item selection within either framework.

\paragraph{Limitations.}
Our analysis has several scope limitations that qualify the generality of findings:

\begin{itemize}[topsep=2pt,itemsep=1pt]
  \item \textbf{Single GPU architecture.} All experiments run on NVIDIA RTX PRO 6000 (Blackwell). Hardware platform effects (A100 vs.\ H100, AMD vs.\ NVIDIA) remain unquantified. Our ``infrastructure'' analysis is limited to numerical precision separability.
  \item \textbf{Four models at 7--9B scale.} Whether variance structure changes with model scale is an open question. Larger models may reduce model$\times$item interaction (convergent capabilities) or increase it (specialized item-level behaviors). Our exchangeability assumption is supported by D-study validation ($\Delta = 0.004$) but remains based on a purposive sample of four architecture families.
  \item \textbf{Four benchmarks, no MATH.} The MATH benchmark was blocked by answer extraction pipeline issues (symbolic expression normalization). Its absence prevents within-free-form format comparison (GSM8K is the sole free-form benchmark), limiting our ability to disentangle format effects from task-content effects.
  \item \textbf{No fine-tuning analysis.} Our analysis addresses inference-time evaluation only. Fine-tuning variance structure---where learning rate, data order, and seed interact with task-specific loss landscapes---is a qualitatively different problem. \citet{bui2025} report extreme seed sensitivity in fine-tuning (variance = 18.22 for RoBERTa on COPA), suggesting that fine-tuning variance structures may be even more complex.
  \item \textbf{Binary scoring primarily.} Logprob analysis is limited to MC formats (MMLU, ARC). Free-form generation (GSM8K) lacks a natural single-token probability score, leaving the scoring-metric dimension partially unresolved.
  \item \textbf{GSM8K estimation instability.} Henderson I convergence issues on GSM8K (negative REML estimates, oscillating components, fixed-vs-random $\Delta = 8.4$ pp) mean that GSM8K variance decomposition should be interpreted as approximate rather than precise.
\end{itemize}

%% ============================================================
%% 7  CONCLUSION
%% ============================================================
\section{Conclusion}

We present a multi-benchmark Generalizability Theory analysis of LLM evaluation variance, analyzing over 800{,}000 evaluation records across four models and four benchmarks.
Three findings emerge.
First, model$\times$item interaction accounts for 33.6\% of variance on MMLU, making item difficulty rankings substantially model-specific and requiring $\sim$10 models for $G \geq 0.90$.
Second, variance structures are irreducibly benchmark-specific---item dominance, model$\times$item magnitude, and required sample sizes follow no systematic pattern across benchmarks, making universal reproducibility recipes illusory.
Third, G-theory requires 4.4$\times$ more items than IRT for reliable measurement because it captures prompt, seed, and ordering dimensions that IRT ignores.
These findings establish that benchmark-specific reproducibility budgets, derived from G-theory D-studies, are the only viable approach to dependable LLM evaluation.
We release an open-source G-theory pipeline and D-study budget calculator to enable practitioners to compute benchmark-specific budgets for their evaluation designs.

%% ============================================================
%% REFERENCES
%% ============================================================
\bibliographystyle{plainnat}
\begin{thebibliography}{25}

\bibitem[Biderman et~al.(2024)]{biderman2024}
Biderman, S., Schoelkopf, H., Anthony, Q., et~al.
\newblock Lessons from the trenches on reproducible evaluation of language models.
\newblock \emph{arXiv preprint}, 2024.

\bibitem[Bouthillier et~al.(2021)]{bouthillier2021}
Bouthillier, X., Delaunay, P., Bronzi, M., et~al.
\newblock Accounting for variance in machine learning benchmarks.
\newblock \emph{Proceedings of Machine Learning and Systems}, 3:747--769, 2021.

\bibitem[Brennan(2001)]{brennan2001}
Brennan, R.~L.
\newblock \emph{Generalizability Theory}.
\newblock Springer, New York, 2001.

\bibitem[Bui et~al.(2025)]{bui2025}
Bui, N., Yu, G., et~al.
\newblock On the sensitivity of fine-tuning to random seeds.
\newblock \emph{arXiv preprint}, 2025.

\bibitem[Dehghani et~al.(2021)]{dehghani2021}
Dehghani, M., Tay, Y., Gritsenko, A.~A., et~al.
\newblock The benchmark lottery.
\newblock \emph{arXiv preprint arXiv:2107.07002}, 2021.

\bibitem[Dodge et~al.(2020)]{dodge2020}
Dodge, J., Ilharco, G., Schwartz, R., et~al.
\newblock Fine-tuning pretrained language models: Weight initializations, data orders, and early stopping.
\newblock \emph{arXiv preprint arXiv:2002.06305}, 2020.

\bibitem[Henderson et~al.(2018)]{henderson2018}
Henderson, P., Islam, R., Bachman, P., et~al.
\newblock Deep reinforcement learning that matters.
\newblock In \emph{Proceedings of AAAI}, 2018.

\bibitem[Hua et~al.(2025)]{hua2025}
Hua, H., Li, Y., et~al.
\newblock Generation-based evaluation of language models.
\newblock In \emph{Proceedings of EMNLP}, 2025.

\bibitem[Jiang \& Skorupski(2017)]{jiang2017}
Jiang, T. and Skorupski, W.~P.
\newblock A {Bayesian} approach to generalizability theory with binary data.
\newblock \emph{Applied Psychological Measurement}, 41(4):287--302, 2017.

\bibitem[Kunievsky \& Evans(2025)]{kunievsky2025}
Kunievsky, D. and Evans, J.
\newblock Decomposing evaluation uncertainty: Intent, articulation, and model components.
\newblock \emph{arXiv preprint}, 2025.

\bibitem[Lalor et~al.(2019)]{lalor2019}
Lalor, J.~P., Wu, H., and Yu, H.
\newblock Learning latent parameters without human response patterns: Item response theory with artificial crowds.
\newblock In \emph{Proceedings of NAACL}, 2019.

\bibitem[Lu et~al.(2022)]{lu2022}
Lu, Y., Bartolo, M., Moore, A., Riedel, S., and Stenetorp, P.
\newblock Fantastically ordered prompts and where to find them: Overcoming few-shot prompt order sensitivity.
\newblock In \emph{Proceedings of ACL}, 2022.

\bibitem[Messina \& Scotta(2026)]{messina2026}
Messina, E. and Scotta, L.
\newblock Floating-point non-associativity in neural network inference.
\newblock \emph{arXiv preprint}, 2026.

\bibitem[Mesting(2026)]{mesting2026}
Mesting, F.
\newblock A generalizability theory analysis of {LLM} evaluation.
\newblock \emph{arXiv preprint}, 2026.

\bibitem[Mushquash \& O'Connor(2006)]{mushquash2006}
Mushquash, C. and O'Connor, B.~P.
\newblock {SPSS} and {SAS} programs for generalizability theory analyses.
\newblock \emph{Behavior Research Methods}, 38(3):542--547, 2006.

\bibitem[Perlitz et~al.(2023)]{perlitz2023}
Perlitz, Y., Bandel, E., Gera, A., et~al.
\newblock Efficient benchmarking of language models.
\newblock \emph{arXiv preprint}, 2023.

\bibitem[Pezeshkpour \& Hruschka(2024)]{pezeshkpour2024}
Pezeshkpour, P. and Hruschka, E.
\newblock Large language models sensitivity to the order of options in multiple-choice questions.
\newblock \emph{arXiv preprint}, 2024.

\bibitem[Polo et~al.(2024)]{polo2024}
Polo, F.~M., Choi, Y., et~al.
\newblock tiny{B}enchmarks: Evaluating {LLM}s with fewer examples.
\newblock \emph{arXiv preprint}, 2024.

\bibitem[Rodriguez et~al.(2021)]{rodriguez2021}
Rodriguez, P., Barrow, J., Hoyle, A.~M., et~al.
\newblock Evaluation examples are not equally informative: How should that change {NLP} leaderboards?
\newblock In \emph{Proceedings of ACL}, 2021.

\bibitem[Sclar et~al.(2024)]{sclar2024}
Sclar, M., Choi, Y., Tsvetkov, Y., and Suhr, A.
\newblock Quantifying language models' sensitivity to spurious features in prompt design.
\newblock In \emph{Proceedings of ICLR}, 2024.

\bibitem[Sellam et~al.(2022)]{sellam2022}
Sellam, T., Yadlowsky, S., Wei, J., et~al.
\newblock Multi{BERT}s: {BERT} reproductions for robustness analysis.
\newblock In \emph{Proceedings of ICLR}, 2022.

\bibitem[Shavelson \& Webb(1991)]{shavelson1991}
Shavelson, R.~J. and Webb, N.~M.
\newblock \emph{Generalizability Theory: A Primer}.
\newblock Sage Publications, 1991.

\bibitem[Shavelson et~al.(1989)]{shavelson1989}
Shavelson, R.~J., Webb, N.~M., and Rowley, G.~L.
\newblock Generalizability theory.
\newblock \emph{American Psychologist}, 44(6):922--932, 1989.

\bibitem[Vania et~al.(2021)]{vania2021}
Vania, C., Kementchedjhieva, Y., S{\o}gaard, A., and Webber, B.
\newblock Comparing test sets with item response theory.
\newblock In \emph{Proceedings of ACL}, 2021.

\bibitem[Yuan et~al.(2025)]{yuan2025}
Yuan, Z., Liu, Y., et~al.
\newblock Hardware matters: Understanding the impact of infrastructure on {LLM} evaluation.
\newblock \emph{arXiv preprint}, 2025.

\end{thebibliography}

\end{document}
